\chapter{Antecedentes y estado del arte}\label{cap:estado}

Este capítulo resume una revisión focalizada y documentada, no una revisión sistemática PRISMA. Las cadenas de búsqueda, las fuentes y los criterios están en el anexo~\ref{anx:busqueda}. Se organiza en cinco frentes, del más cercano al más general:
\begin{enumerate}[label=(\roman*)]
\item los antecedentes peruanos sobre el precio de la vivienda;
\item la cochera como atributo residencial y su contexto en Lima;
\item el paso del aprendizaje automático predictivo al causal;
\item el uso de avisos en línea y de su texto como fuente de datos;
\item la síntesis y el vacío que aborda esta tesis.
\end{enumerate}
Las afirmaciones de originalidad se limitan a las fuentes consultadas.

\section{Marco de referencia: precios hedónicos}

\citet{rosen1974} formalizó la idea de que el precio de un bien diferenciado depende de sus atributos y de que la derivada de esa función respecto de un atributo es su precio implícito. El modelo hedónico se convirtió en la herramienta estándar para valorar atributos de la vivienda \citep{malpezzi2003}. La revisión de \citet{sirmans2005} muestra, sin embargo, que el coeficiente de un mismo atributo varía mucho entre estudios según qué otras variables se controlen. Un precio implícito es \emph{condicional a la especificación}: no indica por sí mismo cuánto cambiaría el precio si se modificara el atributo. Esa distinción entre asociación y efecto organiza todo el capítulo.

\section{Antecedentes peruanos}\label{sec:peru}

\paragraph{Índices hedónicos.} \citet{mundaca2018}, del BCRP, construyen un índice de precios de departamentos \emph{en venta} para diez distritos de Lima (unos 30\,000 avisos, 2008--2016). En sus regresiones semilogarítmicas, el número de cocheras y el de baños son las características con mayor valoración por unidad. Su objetivo es el índice, no identificar efectos. Los propios autores reportan variables omitidas en sus modelos y señalan como principal limitación la falta de un buen indicador de ubicación. \citet{mayo2017} elabora en la Universidad Nacional de Ingeniería un índice hedónico de viviendas \emph{nuevas} con datos del Fondo MiVivienda (2001--2015), y \citet{rivera2018} estudia los determinantes del precio por metro cuadrado (2005--2017). Los tres trabajan con precios de venta y modelos lineales.

\paragraph{Aprendizaje automático predictivo.} \citet{cacnahuaray2024} predicen el alquiler de departamentos en la provincia de Lima con 3\,675 avisos de 2023 y técnicas de AutoML. Entre los predictores más importantes que reportan están el área y el número de cocheras. \citet{barrientos2026} comparan árboles de decisión y XGBoost con la regresión hedónica sobre más de 40\,000 registros del BCRP (2014--2024) e identifican el tamaño, la cochera y la ubicación como los determinantes principales. Ambos confirman que la cochera es relevante para \emph{predecir} el precio en Lima, pero no estiman cuánto lo cambia.

\paragraph{Inferencia causal sobre vivienda en Lima.} El único antecedente local con una estrategia de identificación explícita es \citet{contreras2024}. Con diferencias en diferencias, aprovechan la apertura escalonada de la Línea 1 del Metro para estimar su efecto sobre la informalidad de la vivienda. Su variable de interés no es el precio, pero muestra que el razonamiento causal ya se aplica a la vivienda limeña.

\section{La cochera como atributo residencial}\label{sec:cochera_lit}

\paragraph{Evidencia sobre su valor en alquileres.} La referencia más directa por el tipo de tratamiento es \citet{gabbe2017}. Con más de 38\,000 viviendas en alquiler de áreas metropolitanas de Estados Unidos, estiman que una cochera incluida en el alquiler eleva la renta en alrededor de 17\,\%. \citet{manville2013} encuentra en Los Ángeles que la cochera incluida añade unos 200 dólares mensuales al alquiler. Ambos estudian mercados con una tenencia de automóviles muy superior a la de Lima, por lo que sirven para contrastar la dirección y el orden de magnitud, no para comparar cifras directamente.

\paragraph{Contexto limeño.} Tres hechos locales son relevantes para interpretar el valor de la cochera en Lima:
\begin{itemize}[leftmargin=*]
\item \textbf{Regulación.} En los distritos centrales, la normativa municipal fija un número mínimo de estacionamientos por unidad de vivienda: Miraflores lo hace en el artículo 10 de la Ordenanza 342-MM \citep{miraflores2011} y San Isidro en la Ordenanza 523-MSI \citep{sanisidro2020}. Ahí la cochera tiende a ser un componente estándar de los edificios nuevos.
\item \textbf{Mercado.} El índice de referencia del mercado de alquiler, que también usa el BCRP, toma como unidad típica un departamento de 60 o 100\,m² \emph{con un estacionamiento} \citep{urbania2025}.
\item \textbf{Demanda.} Solo alrededor del 17\,\% de los hogares de Lima y Callao tenía al menos un automóvil en 2015 \citep{limacomovamos2016}, frente a una mayoría en los mercados estadounidenses citados.
\end{itemize}
Estos hechos sugieren que el valor de la cochera en Lima debería ser menor que en Estados Unidos y distinto según la zona. Esta tesis lo pone a prueba.

\section{Del aprendizaje automático predictivo al causal}

\paragraph{Predicción.} Además de los estudios peruanos, la literatura internacional usa el aprendizaje automático para valuar inmuebles. \citet{tapia2025} comparan modelos automáticos de valuación con hedónicos espaciales en Santiago de Chile. \citet{gyger2026} combinan un proceso gaussiano espacial con \textit{tree boosting} para modelar alquileres en Alemania. Estos trabajos mejoran la predicción, pero las importancias de variables o los valores SHAP describen cómo usa el modelo cada variable, no el efecto de intervenirla.

\paragraph{Identificación causal en vivienda.} La economía urbana estima el valor de amenidades con diseños cuasiexperimentales, como discontinuidades en fronteras \citep{black1999,keele2015} o cambios regulatorios \citep{chay2005}. Cuando no hay una fuente de variación casi aleatoria, como ocurre con la cochera de un departamento, la alternativa es ajustar por los confusores observados bajo ignorabilidad condicional y acotar el sesgo por confusores omitidos.

\paragraph{Double Machine Learning.} \citet{chernozhukov2018} muestran que, si se estiman las funciones auxiliares con aprendizaje automático, se ortogonaliza el problema y se usa \textit{cross-fitting}, el efecto es consistente y admite inferencia válida. Para tratamientos binarios, la versión interactiva (IRM/AIPW) es doblemente robusta. El marco tiene limitaciones conocidas. \citet{hunermund2023} muestran que DML es muy sensible a la inclusión de malos controles, y \citet{cinelli2024} sistematizan qué variables deben y no deben ajustarse. \citet{wager2018} y \citet{athey2019} extienden el enfoque a efectos heterogéneos, cuya estimación es frágil en muestras moderadas \citep{knaus2021}.

\paragraph{Robustez y validación.} \citet{cinelli2020} y \citet{chernozhukov2022} formulan cotas de sesgo por confusores omitidos que se calibran con confusores observados. \citet{moran1950} y \citet{cameron2015} proveen herramientas para la dependencia espacial y la inferencia agrupada. \citet{advani2019} recomiendan validar los estimadores con simulaciones ancladas en los datos reales, y \citet{chen2026} lo hacen al comparar regresiones hedónicas y aprendizaje causal en vivienda.

\section{Avisos en línea y su texto como fuente de datos}

\paragraph{Avisos como datos.} \citet{boeing2017} muestran que los avisos de alquiler extraídos por \textit{web scraping} permiten estudiar mercados de alquiler a gran escala. También advierten que exigen limpieza, deduplicación y cuidado al interpretarlos como precios de oferta.

\paragraph{Información en el texto.} \citet{nowak2017} encuentran que el texto libre de los avisos reduce en más de 25\,\% el error de precio de un modelo hedónico, y \citet{shen2021} que la información ``blanda'' de la descripción tiene un valor significativo en el precio. Los atributos relevantes no siempre están en los campos estructurados.

\paragraph{Modelos de lenguaje para medir atributos.} \citet{faxvaag2026} usan modelos GPT para extraer de los avisos de Oslo variables cuantitativas, entre ellas el estacionamiento, e incorporarlas a un modelo de valuación. Es el precedente metodológico más cercano a la medición de la cochera de esta tesis. \citet{egami2023} advierten que usar etiquetas producidas por un modelo de lenguaje directamente en un análisis estadístico produce sesgo e intervalos inválidos, incluso con un 80--90\,\% de acierto. Proponen combinarlas con una muestra etiquetada por personas mediante un estimador doblemente robusto, lo que justifica la validación humana prevista en esta tesis.

\section{Síntesis y vacío}

\begin{table}[H]
\centering\small
\caption{Cuadro comparativo de antecedentes}\label{tab:comparativo}
\begin{tabularx}{\textwidth}{>{\raggedright\arraybackslash}p{2.6cm}>{\raggedright\arraybackslash}p{2.5cm}>{\raggedright\arraybackslash}X>{\raggedright\arraybackslash}p{2.6cm}}
\toprule
Trabajo & Objeto / datos & Qué aporta y límite & Método \\
\midrule
\citet{mundaca2018} & Venta, Lima, 2008--2016 & Cochera y baños, lo más valorado; asociativo, ubicación por distrito & Hedónico \\
\citet{mayo2017}; \citet{rivera2018} & Venta, Lima & Índice y determinantes del precio/m²; asociativo & Hedónico / regresión \\
\citet{cacnahuaray2024} & Alquiler, Lima, 2023 & La cochera, entre los predictores más fuertes; no estima efecto & ML predictivo \\
\citet{barrientos2026} & Venta, Lima (BCRP) & La cochera como determinante clave; predictivo & XGBoost vs. hedónico \\
\citet{contreras2024} & Informalidad, Lima & Único causal local; no estudia precios & Diferencias en diferencias \\
\citet{gabbe2017} & Alquiler, EE.\,UU. & Cochera incluida +17\,\% del alquiler; otro mercado & Hedónico \\
\citet{faxvaag2026} & Avisos, Oslo & Estacionamiento extraído con GPT; objetivo predictivo & LLM + valuación \\
\citet{egami2023} & Metodológico & Uso válido de etiquetas de LLM con validación humana & DSL (doblemente robusto) \\
\textbf{Esta tesis} & \textbf{Alquiler, Lima, 2025} & \textbf{Efecto causal de la cochera medida con IA, con soporte común y sensibilidad} & \textbf{DML-IRM + validación semisintética} \\
\bottomrule
\end{tabularx}
\end{table}

La tabla~\ref{tab:comparativo} resume la posición de la tesis. Los estudios peruanos coinciden en que la cochera es un atributo relevante del precio, pero la abordan con modelos asociativos o predictivos y sobre precios de venta. La literatura internacional ofrece evidencia sobre el valor de la cochera en alquileres y herramientas para estimar efectos causales y medir atributos desde el texto. En las fuentes consultadas no se encontró un estudio sobre Lima que combine las tres piezas: estimar causalmente el valor de un atributo del alquiler, medir ese atributo desde el texto con IA y acotar la fragilidad del resultado. Ese es el espacio en que se sitúa esta tesis.
